\newgeometry{top=2cm, bottom=2cm, left=2cm, right=2cm}
\begin{titlepage}
	\centering

    \rule{0.9\textwidth}{1pt} 

	\vspace{2pt}\vspace{-\baselineskip}

	\rule{0.9\textwidth}{0.4pt}

	\vspace{0.1\textheight}

    \textcolor[rgb]{0.78,0,0}{
    %\textcolor{black}{
    %{\fontfamily{put}\selectfont
    {\Huge REAL-TIME}\\[0.5\baselineskip]
    {\Huge RENDERING}\\[1.7\baselineskip]
    {\Large Realistic Large Bodies of Water}\\
    %{\Large for Large Bodies of Water}\\
    %{\Large Seen at Aerial View}\\
    %{\Large and case study for 0 A.D.}
    }%}

    \vspace{0.025\textheight}
    %\textcolor[rgb]{0.78,0,0}{	
    \rule{0.3\textwidth}{0.4pt}
    %}

    \vspace{2cm}
    \includegraphics[width=0.15\textwidth]{figures/BFH_Logo_BW.pdf}\par\vspace{1cm}
    \vspace{1cm}
    {\Large{\itshape~Samuel Luc Gauthier\/}\par}{}
    \vspace{0.5cm}
    {\scshape Bern University of Applied Sciences\par}
	{\scshape\Large BTI 7302 Project 2\par}

	\vspace{1cm}
	supervised by\par
	Prof.~Marcus \textsc{Hudritsch}
    \vspace{1cm}
	\vfill

    {\large \today\par}%}
    \vfill

    \rule{0.9\textwidth}{0.4pt}

	\vspace{2pt}\vspace{-\baselineskip}

	\rule{0.9\textwidth}{1pt}

\end{titlepage}
\restoregeometry{}

\title{\textbf{Real-Time Rendering}\\\Large{Realistic Large Bodies of Water}}

\author{\normalsize{Samuel Luc Gauthier}}
\affil{Bern University of Applied Sciences \protect\\
\href{mailto:gauts1@bfh.ch}{\nolinkurl{gauts1@bfh.ch}}}
\date{\normalsize{\today}}
\maketitle

\begin{abstract}
The realistic rendering of ocean water constitutes an open problem in computer
graphics. It is a complex topic, as it involves many physical properties. In
this project report we present some solutions from the litterature, and
implement the \textit{Gerstner Wave} model, a deep ocean model, into a
demonstration application. The visual result is not conviencing but provides a
good start and confirms our intuitons. For further work we will implement the
\textit{fast Fourier transform} model and use a derivative of the
\textit{Projected Grid} concept. If there is time left we will handle foam,
caustics and interactions with objects.
\end{abstract}

\small{{\bf Keywords:} computer graphics, real-time simulation, water rendering,
waves, ocean, GPU algorithms}
